\chapter{関連研究}
\revise{本章では，家庭用料理ロボットの操作実行層における本研究の位置づけを明確にするため，先行研究を形状，力学，環境・道具という三つの分析観点と，三つの検証タスクから整理する．まず，料理ロボット全体の研究動向を概観し，レシピ理解やタスク計画と，それらが決定された後の物理操作実行との境界を確認する．次に，各分析観点に関する研究を整理し，最後に皮むき，定量注ぎ，および切断中の安定把持に関連する研究と本論文のタスク固有手法を比較する．}

\section{料理ロボットの研究}
料理ロボットに求められる汎化能力と物理的相互作用の観点について，Sochackiら\cite{Sochacki2024PracticalChefReview}は料理ロボットに関する包括的なレビューを行い，家庭環境での調理では対象物の状態変動と複雑な接触相互作用が支配的であり，単純なピックアンドプレースだけではタスクが成立しないことを指摘している．

調理タスクにおける道具媒介操作と計画の観点では，家庭環境の複雑さに対応するため，Bolliniら\cite{Bollini2013RecipesCookingRobot}やZhangら\cite{Zhang2024RealWorldCookingFromRecipes}は，自然言語のレシピ指示からロボットの行動計画を生成し，統合システム上で実行するアプローチを提案している．
また，Kümpelら\cite{Kumpel2025CookingActions}は調理タスクを特定の操作カテゴリに写像するための知識表現を提案し，Beetzら\cite{Beetz2011CRAM}やPauliusら\cite{PauliusSun2019CookingSurvey}は認知ロボットアーキテクチャや知識表現に基づく計画の重要性を論じている．
さらに，Beckerら\cite{Becker2011KitchenContainers}はキッチン環境でのモバイルマニピュレーションを実装した．
\revise{これらの研究はタスクレベルの計画や論理表現を主眼とするため，計画された個々の操作について，対象物の個体形状や接触条件から軌道，操作量，把持位置，力などを決定する実行層の問題は，別途解く必要がある．}
また近年，Welteら\cite{Welte2025DexterousILSurvey}やWangら\cite{Wang2025EmbodiedVisualPerceptionSurvey}がレビューするように，模倣学習やエンボディドAIに基づく視覚と操作の統合も盛んである．
しかし，料理における道具を介した操作には厳密な物理制約が存在する．
\revise{本論文は，操作目的，使用する道具，および基本条件が決定された後の「動作と制御パラメータの導出」に焦点を当てる．3Dモデルはそのための入力表現であり，形状だけでなく，タスクごとに既知または計測可能な物理条件と組み合わせて用いる．}

\section{幾何学的な分析に基づく操作設計}
家庭環境の多様な形状に適応するための3Dモデル化と環境認識について，ロボットビジョンにおける3D再構成技術は，RGB-Dセンサの発展とともに成熟してきた．
Newcombeら\cite{Newcombe2011KinectFusion}が提案したKinectFusionは，TSDF（Truncated Signed Distance Function）を用いてリアルタイムな密3D再構成を実現した古典的かつ強力な手法である．
その後，Whelanら\cite{Whelan2015ElasticFusion}によるElasticFusionやDaiら\cite{Dai2017BundleFusion}によるBundleFusionにより，非剛体整合やグローバル整合の最適化が進展した．
また，近年の屋内3Dモデリングの動向\cite{Yuan2021RGBDIndoorSurvey}や，計算効率を高めたTSDF融合手法\cite{Li2024RGBTSDF}，さらにはGaussian Splattingを用いたオンラインマッピング\cite{GSFusion2024}やNeRFとRGB-Dを組み合わせた表面再構成\cite{Lee2024InFusionSurf}など，より高品質なモデル化手法が提案されている．
これらの技術は，本研究が前提とする「計算代理としての3Dモデル」を構築するための不可欠な基盤となる．

幾何特徴抽出と軌道生成に関して，取得した点群から，法線や曲率などの幾何特徴を抽出することは，操作設計の第一歩である．
NguyenとLe\cite{Nguyen2013PointCloudSegSurvey}は点群セグメンテーションの包括的なサーベイを提供し，Martínez-Otzetaら\cite{MartinezOtzeta2022RANSACSurvey}はRANSACを用いた平面やモデルの推定の有用性を整理した．
さらに，これらの幾何情報を直接利用した軌道生成として，Yuら\cite{Yu2021PointCloudSlicingRobotica}やZhangら\cite{Zhang2025ContourParallelToolpathPointCloud}は，点群をスライスして断面輪郭を抽出し，ツールパスを生成するアルゴリズムを提案している．
Zengら\cite{Zeng2024ShapeFollowingLaserCleaning}も3D点群表面に基づく形状追跡軌道の計画を行っている．
本論文における「皮むき」タスクの軌道設計は，これらの断面抽出アプローチと軌を一つにするが，さらに力学的接触の維持を統合する点に特徴がある．

\section{力学的な分析に基づく操作設計}
調理操作における対象物の変形や破壊を伴う接触タスクと基本制御について，ロボットが環境と接触する際の制御理論として，Hogan\cite{Hogan1985ImpedanceControl}のインピーダンス制御や，RaibertとCraig\cite{RaibertCraig1981Hybrid}のハイブリッド位置/力制御が広く知られている．
これらの古典的枠組みは，拘束方向と自由方向を分離・管理する上で現在も極めて有効である．

外力が加わる状況下の安定性解析と外乱耐性に関して，切断や保持において外力が加わる場合，把持の安定性評価が必要不可欠である．
Bicchi\cite{Bicchi1995Closure}は力学的なClosure特性を統一的に論じ，FerrariとCanny\cite{FerrariCanny1992OptimalGrasps}はWrench空間に基づく把持の品質指標（Grasp quality measures）を提案した．
RoaとSuárez\cite{RoaSuarez2014GraspQualityReview}はこれらの指標を総括し，外乱に対する抵抗力がタスクの成否を分けることを示している．
本論文の「安定把持」タスクは，まさに切断反力を抵抗する適切な把持点の探索問題である．

食品の切断モデルと把持の近年方向について，食品切断の力学は非常に複雑である．
Mitsioniら\cite{Mitsioni2021FoodCuttingDynamics}は切断ダイナミクスをRNNで学習し制御に組み込んだ．
Muら\cite{Mu2019RoboticCutting}やJamdagniら\cite{Jamdagni2019FractureMechanicsCutting}は，破壊力学に基づき切断過程の摩擦と破断力をモデル化している．
また，Zhouら\cite{Zhou2006PressingSlicingCutting_1, Zhou2006PressingSlicingCutting_2}は，押し込みとスライスを組み合わせた生体材料の切断メカニズムを分析した．
一方，近年の把持研究においては，Songら\cite{Song2025DexterousGraspingOverview}が学習ベースの把持を総覧し，Srinivasら\cite{Srinivas2025GraspFactory}は外乱下での把持ロバスト性を考慮したデータセットを構築している．
本論文は，これらの切断力学モデルを外力として取り込み，把持品質の解析的評価関数に組み込むことで，安定した切断中の把持を実現する．

\section{環境分析に基づく操作設計}
環境拘束の積極的な利用とタスク・モーションプランニングに関して，Garrettら\cite{Garrett2021TAMPSurvey}やZhaoら\cite{Zhao2024OptTAMPSurvey}のレビューが示す通り，離散タスクと連続運動の統合計画において，幾何的制約と衝突回避は中核的な課題である．
Stilman\cite{Stilman2011ConstrainedManipulation}は環境拘束下での操作計画を定式化し，Görnerら\cite{Gorner2019MoveItTaskConstructor}はこれらを実装可能なフレームワークとして提供している．
さらに，Nakatsuruら\cite{Nakatsuru2023EnvironmentSupport}は，壁や机などの環境支持を積極的に利用する計画手法を提案し，
Ard{\'o}nら\cite{Ardon2021AffordanceReview}も環境と対象の相互作用可能性を表現するアフォーダンス関係の学習・表現についても，設計選択と汎化の観点から体系的な整理が進められている．

不確実性の高いキッチン環境における能動知覚と環境モデリングについて，Bohgら\cite{Bohg2016InteractivePerception}が提唱するインタラクティブ知覚のように，ロボット自身の行動を通じて観測を改善するアプローチが重要となる．
本論文におけるモデルの計測方針は，対象周囲の複数視点からの点群取得に加え，必要に応じて対象を把持して底面側を観測し，取得した点群を統合することで不完全観測を低減するという実践的なアプローチとして位置づけられる．

\section{タスクごとの先行研究}
本節では，本論文で検証タスクとして取り上げる「注ぐ」「皮むき」「安定把持」の各操作に関する個別研究を整理し，提案手法の有効性を明らかにする．

\subsection{皮むきに関する研究}
皮むき動作は，不規則な曲面に対する幾何学的な軌道追跡と，接触力を維持する力制御の統合が必要なタスクである．
Houら\cite{Hou2022PeelingPlanning}は自動皮むきのための幾何的な経路設計問題を整理した．
また，マルチモダルな情報とインピーダンス制御を用いた単腕皮むきシステム\cite{Arxiv2024MultimodalPeeling}や，対象物の再配向を伴う拘束付き巧緻操作としての皮むき\cite{Arxiv2024ConstrainedDexterousPeeling}も実機実装として報告されている．

剥皮を成立させる前提となる手内操作（In-hand manipulation）能力については，Rojasら\cite{Rojas2016GR2}の劣駆動機構による手内平面操作，Maら\cite{Ma2017WholeHandCaging}のCagingによるロバスト把持，McCannとDollar\cite{McCann2017StewartHand}のパラレルメカニズム型ハンド，Mnyusiwallaら\cite{Mnyusiwalla2016BioInspiredInsideHand}の生体模倣指などが提案されている．
さらにLiarokapisとDollarは，手内操作のプリミティブ抽出\cite{Liarokapis2017Primitives}やタスク特化モデルの学習\cite{Liarokapis2016TaskSpecificModels}を論じている．
微小力センシングの観点では，網膜手術などの医療分野における微小力フィードバック研究\cite{Balicki2010MicroForcePeeling}が，接触力管理の設計思想として参考になる．
\revise{本論文の皮むき操作は，食材ごとの形状から到達可能な軌道を生成し，表面法線に基づく領域分割と既処理境界に基づく食材回転を行い，さらに力覚フィードバックで点群誤差と表面凹凸を補償する点に特徴がある．}

\subsection{注ぐに関する研究}
ロボットによる注ぎ操作は，液体モデルの複雑さと容器の多様性から多くの研究がなされてきた．
Rozoら\cite{Rozo2013PouringForceBased}は力覚に基づく模倣学習を，Chenら\cite{Chen2019RNNMPCPouring}はRNNとMPCを組み合わせた注ぎ制御を提案した．
また，López-Guevaraら\cite{LopezGuevara2017NotToSpill}は近似流体モデルを用いた強化学習を行い，衝突回避を含めた注ぎの動作計画\cite{AAAI2023PouringPlanning}も提案されている．
Reyes-Montiel\cite{ReyesMontiel2026PouringReview}の最新のレビューやHuangら\cite{Huang2021SelfSupervisedPouring}の自己教師あり学習の事例は，流体操作における学習ベース手法の隆盛を示している．

一方，解析的なアプローチとして，DoとBurgard\cite{DoBurgard2018AccuratePouring}やSchenckとFox\cite{SchenckFox2017ClosedLoopPouring}，Kennedyら\cite{Kennedy2017Dispensing}は，RGB-Dカメラや視覚フィードバックを用いた高精度な閉ループ制御を報告している．
容器の幾何形状に着目した研究としては，Brandlら\cite{Brandl2014WarpedPouring}のワーピングによる一般化や，Mottaghiら\cite{Mottaghi2017GlassHalfFull}の画像からの容積推定がある．
液面知覚に関しては，ステレオ視による流動観測\cite{YamaguchiAtkeson2016StereoFlow}，深層学習による液体知覚\cite{SchenckFox2016LiquidsISER}，RGB-Dを用いた確率的液面推定\cite{Do2016ProbLiquidLevel}などが提案されており，失敗からの回復を学習に組み込む試み\cite{Langsfeld2014FailureToSuccessPouring}も存在する．
\revise{本論文の注ぎ操作は，事前登録した容器形状や制御用の外部計量器具に依存せず，対象容器の3D点群から液面高さと体積の関係，または容器角度と流出量の関係を導出し，指定量に対応する制御目標へ変換する点に特徴がある．}

\subsection{切断時の安定把持に関する研究}
切断における把持安定化の研究は，切断自体の制御と把持計画の融合領域に位置する．
切断制御の古典としてZengとHemami\cite{ZengHemami1997Cutting}やJungとHsia\cite{JungHsia1999Cutting}による適応制御があり，近年では視覚ベース切断の課題整理\cite{Singh2025VisionBasedCuttingSurvey}や，前述の学習動力学\cite{Mitsioni2021FoodCuttingDynamics}が提案されている．
切断力学に基づく刃物の運動最適化については，MuとJia\cite{MuJia2022PropertyKnifeOpt}の物性推定や，Zhouらの詳細な応力分布・力学モデル\cite{Zhou2006PressSliceI, Zhou2006PressSliceII}，Atkinsら\cite{Atkins2004CheddarSalami}の柔軟材料における押し引き切断の分析が重要である．
視覚と力覚を統合した柔軟物切断制御\cite{Long2014ForceVisionCutting}も存在する．

外力に対する安定把持計画の基礎としては，Millerら\cite{Miller2003ShapePrimitives}の形状プリミティブ近似，MirtichとCanny\cite{MirtichCanny1994OptGrasps}の実用的最適把持計算，Pelossofら\cite{Pelossof2004SVMGrasping}のSVMを用いた学習把持，Nguyen\cite{Nguyen1988ForceClosure}のForce-Closureの幾何的構成法が挙げられる．
多指ハンドにおける力の分離については，YoshikawaとNagai\cite{YoshikawaNagai1987Forces}の操作力と把持力の概念や，Donnerら\cite{Donner2018WrenchDecomposition}の物理的に妥当なWrench分解手法が強力な理論的背景となる．
さらに，拘束環境下での衝突を考慮した把持\cite{Lou2021CollisionAwareGrasping}や，重心と力覚に基づく姿勢適応\cite{Kanoulas2018CoMGraspAdapt}など，環境制約と外乱を同時に扱う技術が進展している．
\revise{本論文では，形状から衝突しない把持点を選ぶだけでなく，既知の包丁軌道に沿って時間変化する切断Wrenchを推定し，摩擦および力・モーメントの釣合いとともに把持安定性評価へ組み込む点に特徴がある．}

\section{まとめ}
以上の関連研究は，大きく次の三つのアプローチとして整理できる．
第一に，RGB-D等のセンシングに基づく3Dモデリングと幾何学的解析を基盤とし，対象形状の推定や表面特徴量の抽出を通じて，軌道生成や作業パラメータの導出へ接続する方向性である
\cite{Yuan2021RGBDIndoorSurvey,Li2024RGBTSDF,GSFusion2024,Lee2024InFusionSurf,Nguyen2013PointCloudSegSurvey,MartinezOtzeta2022RANSACSurvey,Yu2021PointCloudSlicingRobotica,Zhang2025ContourParallelToolpathPointCloud,Zeng2024ShapeFollowingLaserCleaning}．
第二に，接触を伴う操作を安定化するための力制御・インピーダンス制御，および切断や把持に関する力学モデル・評価指標を整備し，外力や接触状態の変動に対してロバストな実行を目指す方向性である
\cite{Hogan1985ImpedanceControl,RaibertCraig1981Hybrid,Bicchi1995Closure,FerrariCanny1992OptimalGrasps,RoaSuarez2014GraspQualityReview,Mu2019RoboticCutting,Jamdagni2019FractureMechanicsCutting,Zhou2006PressingSlicingCutting,Long2014ForceVisionCutting,Mitsioni2021FoodCuttingDynamics,MuJia2022PropertyKnifeOpt}．
第三に，道具や作業台を含む環境との相互作用を前提とし，拘束条件の明示化，能動的な再計測，および相互作用可能性（アフォーダンス）の表現を通じて，ツール媒介操作を成立させる方向性である
\cite{Bohg2016InteractivePerception,Ardon2021AffordanceReview,Stilman2011ConstrainedManipulation,Nakatsuru2023EnvironmentSupport,Mottaghi2017GlassHalfFull}．

\begin{revisionblock}
これらの研究により，形状認識，接触制御，環境利用に関する知見はそれぞれ蓄積されてきた．
一方，皮むき，注ぎ，切断中の把持では，3Dモデルから求める量も，最終的な動作・制御パラメータも異なるため，一般的な「計測，分析，動作生成」の流れだけでは各操作は成立しない．
本論文の位置づけは，形状，力学，環境・道具という共通の観点で相補的な事例を選び，皮むきでは表面追従と力制御，注ぎでは容器幾何と操作量，安定把持では切断Wrenchと接触力という，タスク固有の対応関係を構成して実機で検証する点にある．
次章では，これらに共通して必要となるロボットシステムとRGB-D点群の取得・処理・評価方法を示す．
\end{revisionblock}








